\section{Intellect and Knowledge}\label{sec:knowledge}

Having treated what belongs to the angel's substance, Aquinas turns to
the angels' knowledge, \latin{procedendum est ad cognitionem ipsius}, and
divides the treatise fourfold: what belongs to the power of knowledge
(\latin{virtus cognoscitiva}), what belongs to the medium of knowing
(\latin{medium cognoscendi}), the things known, and the mode of their
knowledge (\latin{de modo cognitionis ipsorum}) (\ST{I q.~54 pr.}). The
order is deliberate. Because the angel is a subsisting form with no matter
(\S\ref{sec:q50}), its understanding cannot be an act of a bodily organ,
cannot begin from sense, and cannot be parcelled into the faculties of the
human soul; the five questions that follow work out what such a knowledge
is, first in its power (q.~54), then in the \latin{species intelligibiles}
that inform it (q.~55), then in its objects, immaterial (q.~56) and
material, future, hidden, and gracious (q.~57), and finally in the manner
of its act, which the tradition, after Augustine, names morning and
evening knowledge (q.~58). The Catechism states the ground of the whole
tract in one sentence: ``As purely spiritual creatures angels have
intelligence and will: they are personal and immortal creatures''
(CCC~330).

\subsection{The Angelic Power of Knowledge (\ST{I q.~54})}\label{sec:q54}

The five articles ask whether the angel's act of understanding is his
substance; whether his being is his understanding; whether his substance
is his power of intelligence; whether there is in him an active and a
passive intellect; and whether he has any power of knowledge besides the
intellect (\ST{I q.~54 pr.}).

\paragraph{Understanding is not the angel's substance (a.~1).} Aquinas
determines that ``it is impossible for the action of an angel, or of any
creature, to be its own substance'' (\latin{impossibile est quod actio
Angeli, vel cuiuscumque alterius creaturae, sit eius substantia}). The
governing reason is the metaphysics of act and potency established in
q.~50: an action is the actuality of a power as existence is the actuality
of an essence, and nothing with an admixture of potentiality can be its
own actuality; God alone is pure act, so that in God alone are substance,
existence, and action the same. A second argument shows what the contrary
would cost: a subsisting act of understanding could be but one, so that an
angel identical with his own understanding would be identical with God and
with every other angel, and there could be no degrees of more and less
perfect understanding. The replies free the authorities. When the active
intellect of the soul is said to be its own action, ``such predication is
not essential, but concomitant,'' because its nature stands in act
(ad~1). When Aristotle
calls the intellect's act life, and when Augustine joins memory,
understanding, and will as ``one essence, one life,'' the name of life is
transferred to the essence, as a race stands to running, not identified
with it (ad~2; Aquinas cites \emph{De Trinitate} X). An immanent action is
not a real medium between agent and object but follows upon their union
(ad~3).

\paragraph{In the angel, to understand is not to exist (a.~2).} The
\emph{sed contra} is Dionysius, who speaks of the divine minds as moved:
``the divine minds are said to be moved circularly indeed, by being united
to the illuminations of the Beautiful and Good, without beginning and
without end'' (\emph{DN} 4.8); but to exist is not movement. Aquinas
distinguishes the twofold class of action, transient and immanent.
Immanent actions such as understanding and willing bear, of themselves,
upon all things, because the true and the good are convertible with being;
but ``the being of every creature is restricted to one in genus and
species; God's being alone is simply infinite.'' Hence ``the Divine nature
alone is its own act of understanding and its own act of will.'' The
angelic essence is the principle of the angel's entire existence, but of
his understanding only as included under the universal object of truth
(ad~2).

\paragraph{The power of intelligence is not the angel's essence (a.~3).}
The \emph{sed contra} is Dionysius, \emph{CH} 11, where Aquinas reads that
the angels ``are divided into substance, power, and operation''
(\latin{dividuntur in substantiam, virtutem et operationem}); Parker's
rendering has the Divine Minds ``distributed into three,----into essence,
and power, and energy'' (\emph{CH} 11.2). The reason is universal: since
every power is ordained to its act, diversity of acts means diversity of
powers; in every creature the essence differs from existence as potency
from act; and since in the angel to understand is not to exist (a.~2),
``the angel's essence is not his power of intelligence: nor is the essence
of any creature its power of operation.'' Dionysius's habit of calling the
angels ``intellects'' and ``minds'' is explained, not corrected: an angel
is so called ``because all his knowledge is intellectual,'' whereas the
soul's knowledge is partly sensitive (ad~1). A created simple form, not
being its own existence, can be the subject of the accidents that follow
its species (ad~2). Damascene's definition stands with the determination:
``An angel, then, is an intelligent essence, in perpetual motion, with
free-will, incorporeal, ministering to God'' (\emph{De fide orthodoxa}
II.3) --- essence one thing, its perpetual intellectual motion another.

\paragraph{No active and passive intellect in the angels (a.~4).} The two
intellects are postulated in us for two reasons: we understand some things
only potentially, which requires a passive intellect, and the natures we
understand exist outside the soul only potentially intelligible, which
requires an active intellect to render them actual. ``Each of these
necessities is absent from the angels.'' What they naturally apprehend
they apprehend without potency, and their objects, being first and
principally immaterial things, are actually intelligible; therefore ``there
cannot be an active and a passive intellect in them, except equivocally.''
The objection from angelic illumination is answered precisely: the active
intellect in us enlightens not another intellect but potentially
intelligible things, by abstraction; one angel's enlightening another
(\S\ref{sec:illumination}) is no more that than is the angels' potency to
supernatural revelation the passive intellect (ad~2). Damascene supplies
the received picture Aquinas is refining: the angels ``are secondary
intelligent lights derived from that first light which is without
beginning, for they have the power of illumination'' (\emph{De fide
orthodoxa} II.3).

\paragraph{No knowledge in the angels but intellectual (a.~5).} The
\emph{sed contra} is Gregory: ``man senses in common with the brutes, and
understands with the angels.'' The homily itself argues that man
has something in common with every creature: \latin{commune esse cum
lapidibus, vivere cum arboribus, sentire cum animalibus, intelligere cum
angelis} --- being with stones, life with the trees, sensation with the
animals, understanding with the angels (Gregory, \emph{Hom.\ in Ev.}\
29.2; the same sentence says of the angels themselves that they
\latin{sunt, vivunt, sentiunt et discernunt}, are, live, sense, and
discern). Aquinas's determination: some powers of the soul operate through
corporeal organs, as sight through the eye; others, intellect and will,
through none. ``Now the angels have no bodies naturally joined to them''
(\S\ref{sec:bodies}); hence ``of the soul's powers only intellect and will
can belong to them'' (\latin{de viribus animae non possunt eis competere
nisi intellectus et voluntas}). The order of the universe itself is served
when the highest intellectual creature is entirely intelligent. The
contrary authorities are read two ways: either as following the opinion
that angels and demons have bodies naturally united, which Augustine
entertains without asserting, or as similitudes --- ``sentence'' and
``sensing'' the truth, the ``experience'' of singulars known without
sense, memory as Augustine places it in the mind. Dionysius's ``perverted
phantasy'' in the demons (\emph{DN} 4) is such a similitude: Parker
renders the place, ``what is evil in demons? An irrational anger----a
senseless desire----a headlong fancy'' (\emph{DN} 4.23); the demons'
deception is a false practical estimate of the true good, not the work of
an imagination they do not have.

\angeldossier{\ST{I q.~54 aa.~1--5}; parallels \emph{Summa contra
gentiles} II.96--97; \emph{De veritate} q.~8.}
 {Church teaching: the angels are purely spiritual creatures having
 intelligence and will, personal and immortal (CCC~330). Thomist
 position: in the angel substance, power, and act of understanding are
 really distinct (aa.~1--3); there is no agent and possible intellect in
 them (a.~4); their only cognitive power is the intellect (a.~5).}
 {Dionysius, \emph{DN} 4.8 (\emph{sed contra} of a.~2), \emph{CH} 11.2
 (\emph{sed contra} of a.~3), \emph{DN} 4.23 (obj.~3 of a.~5); Gregory,
 \emph{Hom.\ in Ev.}\ 29.2 (\emph{sed contra} of a.~5); Damascene,
 \emph{De fide orthodoxa} II.3; Augustine, \emph{De civitate Dei} VIII and
 XXI and \emph{De Trinitate} X as cited by Aquinas (aa.~1, 5); Boethius
 and Aristotle as reported by Aquinas; CCC~330.}
 {Averroes and the ancients who made the active intellect its own action
 (obj.~1 of a.~1); the opinion, used but not asserted by Augustine, that
 angels and demons have bodies naturally united to them (a.~5).}

\subsection{The Medium of Angelic Knowledge (\ST{I q.~55})}\label{sec:q55}

Three articles: whether the angels know all things by their substance or
by species; if by species, whether by connatural species or by species
drawn from things; and whether the higher angels know by more universal
species than the lower (\ST{I q.~55 pr.}).

\paragraph{Not all things by their substance (a.~1).} The \emph{sed
contra} is Dionysius, that the angels ``are enlightened by the forms of
things'' (\emph{DN} 4). The medium by which an intellect understands is
its form, and a faculty is perfected by its form only if all the faculty
extends to is contained under that form. The angelic intellect extends to
universal being; but the angel's essence ``is an essence restricted to a
genus and species,'' comprising all things only under a common formality.
``Therefore God alone knows all things by His essence. But an angel cannot
know all things by his essence; and his intellect must be perfected by
some species in order to know things.'' Dionysius's own sentence, that the
angels know things ``according to the proper nature of a mind,'' names the
knowing power, not the medium, which is always a likeness of the thing
known (ad~1). The Aristotelian identity of intellect and thing understood
in what is without matter means that the intellect in act is its object in
act, the species standing as its form (ad~2). When Dionysius says God
``enfolds the whole in the whole'' --- Parker's text has the Cause ``which
has embraced the whole in one'' (\emph{DN} 4.7) --- the angelic essence
holds lower things in a higher and higher things in a participative mode,
but only God's essence holds all things perfectly and according to their
own formality (ad~3).

\paragraph{Species connatural, not drawn from things (a.~2).} The
\emph{sed contra} is Dionysius, that the angels ``do not gather their
Divine knowledge from things divisible or sensible''; in Parker's English
the angelic minds collect their divine knowledge
``not in portions, or from portions, or sensible perceptions, or detailed
reasonings'' (\emph{DN} 7.2). Aquinas's reason mounts by the ladder of
perfection: as the highest bodies have their potentiality fully perfected
by their form while lower bodies receive now one form, now another, so the
higher intellectual substances have their power of understanding
``naturally complete by intelligible species'' --- \latin{species
intelligibiles connaturales}, connatural to them --- while human souls are
completed successively by species drawn from things. The determination:
``The species whereby the angels understand are not drawn from things, but
are connatural to them'' (\latin{species per quas Angeli intelligunt, non
sunt a rebus acceptae, sed eis connaturales}). The same follows from their
mode of existence: souls, as forms of bodies, fittingly seek their
intelligible perfection from bodies; the angels ``are utterly free from
bodies, and subsist immaterially and in their own intelligible nature,''
and receive the species of things
``through an intelligible outpouring'' from God together with their
intellectual nature (\latin{per intelligibilem effluxum, quo a Deo species
rerum cognitarum acceperunt}). Augustine himself states
the order of creation this implies: \latin{caetera vero quae infra sunt
ita creantur, ut prius fiant in cognitione rationalis creaturae, ac deinde
in genere suo} --- the things below the angels are so created that they
exist first in the knowledge of the rational creature and then in their
own kind (\emph{De Genesi ad litteram} II.8.16). The likenesses in the
angelic mind are not images caused by creatures but come ``from God, Who
is the cause of creatures, and in Whom the likenesses of creatures first
exist''; so Augustine again: \latin{ratio qua creatura conditur, prior est
in Verbo Dei quam ipsa creatura quae conditur: sic et eiusdem rationis
cognitio prius fit in creatura intellectuali} (\emph{De Genesi ad
litteram} II.8.17) (ad~1). The angel could not abstract species from
material things even if he would, having no imagination; and he would not,
having no need (ad~2). His knowledge is indifferent to near and distant;
he is moved locally not to acquire knowledge but to operate (ad~3).

\paragraph{The higher angels know by fewer and more universal species
(a.~3).} The \emph{sed contra} joins Dionysius and the \emph{Liber de
causis}: the higher angels have a more universal knowledge and more
universal forms. Parker's \emph{CH} 12.2 gives the substance: the Cherubim
participate in higher wisdom and knowledge, while ``the Divisions of the
Beings beneath them, participate, they also, in wisdom and knowledge, but
nevertheless partially \dots\ in a lower degree.'' The governing reason is
nearness to God. In God the whole plenitude of knowledge is contained in
one thing, the Divine essence; created intellects possess that plenitude
``in a lower manner, and less simply,'' so that a lower intellect needs
many forms where a higher needs few. ``Thus the higher the angel is, by so
much the fewer species will he be able to apprehend the whole mass of
intelligible objects.'' Even among men ``there are others of stronger
intellect, who can grasp many things from few.'' The replies secure the
account against misunderstanding: universality here is not abstracted from
particulars but pre-exists them, causally in the Word, by nature in the
angelic mind (ad~1); to know by a universal \emph{medium} is more perfect,
not less, for the intellect that knows each thing through one universal
medium surpasses one that cannot (ad~2); and an eminent species can be the
proper type of many, as the Divine essence is the proper type of each
thing (ad~3). \emph{De veritate} q.~8 a.~10 argues the same from the
principle that a higher power works through fewer principles that reach
further.

\angeldossier{\ST{I q.~55 aa.~1--3}; parallels \emph{De veritate} q.~8
aa.~8--10; \emph{Summa contra gentiles} II.96, 98, 101.}
 {Thomist position: the angels know by connatural, infused
 \latin{species intelligibiles}, not by their essence alone (a.~1) nor by
 species drawn from things (a.~2); the higher angels know through fewer
 and more universal species (a.~3).}
 {Dionysius, \emph{DN} 4 (\emph{sed contra} and obj.~3 of a.~1), \emph{DN}
 7 (obj.~1 and ad~1 of a.~1; \emph{sed contra} of a.~2), \emph{CH} 12
 (\emph{sed contra} of a.~3); \emph{Liber de causis} (a.~3); Augustine,
 \emph{De Genesi ad litteram} II.8.16--17 (a.~2).}
 {Bonaventure agrees that the angels know all things by innate species,
 yet holds that the angel could receive new species were God to create a
 new nature, and that singulars are known by composing innate species
 while the angel directs its gaze upon the thing (\emph{In II Sent.} d.~3
 p.~2 a.~2 q.~1); Aquinas's account excludes both composition and
 reception (aa.~2--3; \ST{I q.~58 a.~4}).}

\subsection{The Angel's Knowledge of Immaterial Things (\ST{I
q.~56})}\label{sec:q56}

Three articles: whether an angel knows himself; whether one angel knows
another; whether the angel knows God by his own natural principles
(\ST{I q.~56 pr.}).

\paragraph{The angel knows himself (a.~1).} The \emph{sed contra} is
Augustine, that ``the angel knew himself when he was established, that is,
enlightened by truth''; of the light first created he writes,
\latin{in ipsa sua conformatione cognovit, hoc est illustratione
veritatis} --- in its very conformation it knew itself, that is, in the
illumination of the truth (\emph{De Genesi ad litteram} II.8.16). The
determination follows from the metaphysics of immanent action: for such
action the object must be united to the agent as its form, and a
subsisting intelligible form needs no reception to be actually
intelligible. ``And since an angel is immaterial, he is a subsisting form;
and, consequently, he is actually intelligible. Hence it follows that he
understands himself by his form, which is his substance'' (\latin{per suam
formam, quae est sua substantia, seipsum intelligat}). The first reply
carries a textual notice: where the old Latin translation of Dionysius
read that the angels ``do not know their own powers,'' the amended
translation reads that they ``knew their own powers''; Parker's English
follows the corrected text, ``they know their own proper powers and
illuminations'' (\emph{CH} 6.1). Even the older letter could stand, as
meaning that the angels do not comprehend the order of the Divine Wisdom
from which their power proceeds (ad~1). Singularity is no bar: we fail to
know corporeal singulars because of matter, not because of particularity
(ad~2); and the intellect is moved by its object only so far as it is in
potency, which does not touch the angel's knowledge of himself (ad~3).

\paragraph{One angel knows another (a.~2).} The \emph{sed contra} is the
\emph{Liber de causis}: ``every intelligence knows the things which are
not corrupted.'' The determination is Augustinian in frame. What
pre-existed eternally in the Word came forth in two ways: into the angelic
mind, and to subsist in its own nature. God impressed on every angelic
mind the forms of all He produced, corporeal and spiritual --- in each
angel the form of his own species ``according to both its natural and its
intelligible condition,'' so that he subsists in his species and knows
himself by it, and the forms of all other natures according to their
intelligible condition alone. One angel therefore knows another by that
angel's species existing in his intellect, which differs from the angel
himself ``not according to material and immaterial nature, but according
to natural and intentional existence'' (\latin{esse intentionale}), as the
color on the wall has natural existence and in the medium only intentional
(ad~3). The angelic natures, differing only in degrees of intellectual
perfection, do not bar one another from knowledge (ad~1); likeness without
causality suffices (ad~2). And the universe's proportion answers the
cleverest objection: had God resolved to create more angels, He would have
impressed more species, ``as a builder who, if he had intended to build a
larger house, would have made larger foundations'' (ad~4).

\paragraph{The angel knows God by natural principles, but as in a mirror
(a.~3).} The \emph{sed contra} argues from the lesser to the greater: men
know God by natural principles (Rom 1:19), and ``the angels are mightier
in knowledge than men.'' A thing is known in three ways: by the presence
of its own essence in the knower, as God knows Himself and as the blessed
see God; by the presence of the thing's likeness in the knowing power; and
by a likeness reflected in something else, as a man is seen in a mirror.
The angel's natural knowledge of God stands midway. ``Since God's image is
impressed on the very nature of the angel in his essence, the angel knows
God in as much as he is the image of God. Yet he does not behold God's
essence; because no created likeness is sufficient to represent the Divine
essence.'' His knowledge ``approaches rather to the specular kind; because
the angelic nature is itself a kind of mirror representing the Divine
image'' (\latin{ipsa natura angelica est quoddam speculum divinam
similitudinem repraesentans}). The objections are Dionysius's own
superlatives: God is ``altogether incomprehensible to all, and of It,
there is neither perception nor imagination, nor surmise, nor name, nor
expression, nor contact, nor science'' (\emph{DN} 1.2) --- which Aquinas
reads of the knowledge of comprehension, excluded for every created
intellect (ad~1). Paul's ``through a glass in a dark manner'' (1 Cor
13:12) divides vision in the Word from vision in creatures; the angel's
natural knowledge is likened to the specular, though it is not gathered
from sensible things (ad~3). Damascene says the same in the mode of the
Greek tradition: ``They behold God according to their capacity, and this
is their food'' (\emph{De fide orthodoxa} II.3).

\angeldossier{\ST{I q.~56 aa.~1--3}; parallels \emph{De veritate} q.~8
aa.~1--7; \emph{Summa contra gentiles} II.98.}
 {Common teaching: the angel knows himself and the other angels
 (aa.~1--2). Church teaching: no created intellect comprehends God or
 sees his essence by nature (a.~3; \ST{I q.~12}). Thomist position: the
 angel's natural knowledge of God is specular, the angelic nature being a
 mirror of the divine image (a.~3).}
 {Augustine, \emph{De Genesi ad litteram} II.8.16 (\emph{sed contra} of
 a.~1); Dionysius, \emph{CH} 6.1 (obj.~1 of a.~1), \emph{DN} 1.2 (obj.~1
 of a.~3), \emph{DN} 7 (a.~3); \emph{Liber de causis} (\emph{sed contra}
 of a.~2, obj.~2); Rom 1:19--20, 1 Cor 13:12; Damascene, \emph{De fide
 orthodoxa} II.3.}
 {The old Latin version of \emph{CH} 6, ``they do not know their own
 powers,'' corrected in the newer translation (a.~1 ad~1).}

\subsection{The Angel's Knowledge of Material Things (\ST{I
q.~57})}\label{sec:q57}

Five articles: whether the angels know the natures of material things;
whether they know singulars; whether they know the future; whether they
know the thoughts of hearts; whether they know all mysteries of grace
(\ST{I q.~57 pr.}).

\paragraph{The angels know material things (a.~1).} The established order
of things is that the higher contain, eminently and more simply, what the
lower contain deficiently and in manifold manner. In God all things
pre-exist supersubstantially; ``among other creatures the angels are
nearest to God,'' and Dionysius testifies that the holy orders ``share in
the supremely Divine participation, in a higher degree'' than what merely
exists or lives (\emph{CH} 4.2). ``Consequently, all material things
pre-exist in the angels more simply and less materially even than in
themselves, yet in a more manifold manner and less perfectly than in
God.'' Since whatever is in a subject is there after the subject's manner,
and the angels are intellectual by nature, they know material things
forasmuch as these are in them by intelligible species. The replies cut
away the contrary picture of knowledge: the thing understood perfects the
knower precisely by its species (ad~1); the intellect alone apprehends the
essences of things, the \latin{quod quid est}, about which it does not err
(ad~2); and no abstraction is required, because the angel does not draw
his knowledge from material things but has it by connatural species
actually intelligible (ad~3; q.~55 a.~2).

\paragraph{The angels know singulars (a.~2).} The \emph{sed contra} is the
guardianship: ``No one can guard what he does not know,'' and the angels
guard individual men --- ``For he hath given his angels charge over thee;
to keep thee in all thy ways'' (Ps 90[91]:11, Douay). Aquinas first names
the positions he rejects. To deny the angels all knowledge of singulars
``derogates from the Catholic faith, which asserts that these lower things
are administered by angels'' (Heb 1:14), and against the word of
Ecclesiastes, ``say not before the angel: There is no providence''
(Eccles 5:5, Douay); it is also against the philosophers, who made the
separate substances movers of the spheres. Nor does it suffice to know
singulars in their universal causes --- the opinion \emph{De veritate}
q.~8 a.~11 assigns to Avicenna --- as the astronomer foretells an eclipse:
such knowledge does not reach the singular as it exists here and now, and
``administration, providence and movement are of singulars, as they are
here and now existing.'' The determination: as man by diverse powers knows
universals and singulars, so the angel knows both by his one intellective
power; for ``the higher a thing is, so much the more is its power united
and far-reaching,'' as the common sense in man knows all that the five
outer senses know. Hence ``it is unreasonable to say that a man knows by
any one of his powers something which an angel by his one faculty of
knowledge, namely, the intellect, does not know'' --- Aristotle ``pronounces
it ridiculous to say that a discord, which is known to us, should be
unknown to God.'' The manner: God is cause of the entire substance of
things, matter and form, and His knowledge is their cause; the angelic
species, ``the manifold representations of that one simple essence,'' are
images of things not only in their universal natures but in their
individuating principles (ad~2).

\paragraph{The angels do not know the future in itself (a.~3).} The
\emph{sed contra} is Isaiah: ``Shew the things that are to come hereafter,
and we shall know that ye are gods'' (Isa 41:23, Douay); to know future
events is the exclusive sign of the Divinity. The future is known in two
ways. In its causes: what proceeds necessarily from its causes is known
with certainty, as that the sun will rise tomorrow; what follows for the
most part is known conjecturally, as the physician foresees the course of
a disease --- and the angels know the future in this way by so much the
more as they grasp the causes of things more universally and perfectly;
chance events are quite unknown. In themselves (\latin{futura in
seipsis}): ``To know the future in this way belongs to God alone \dots\
for God sees all things in His eternity, which, being simple, is present
to all time, and embraces all time'' --- \latin{solius Dei est futura
cognoscere}. ``But the mind of an angel, and every created intellect, fall
far short of God's eternity.'' Hence the Lord's word stands without
exception: ``But of that day and hour no one knoweth: no, not the angels
of heaven, but the Father alone'' (Matt 24:36, Douay). The replies
discipline the imagination of an angelic eternity: though the angelic
intellect is above the time that measures corporeal motion, there is
succession in its concepts, for as Augustine says, God ``moves the
spiritual creature according to time'' --- in his own words, \latin{per
tempus movet conditum spiritum ipse nec per tempus nec per locum motus
conditor Spiritus} (\emph{De Genesi ad litteram} VIII.20.39) (ad~2).
Innate species are indifferent to present, past, and future, but future
things have not yet the nature that would liken them to the species (ad~3);
distance in place belongs to what already exists (ad~4). \emph{De malo}
q.~16 a.~7 carries the same division for the demons' foreknowledge, with
the same Isaian seal.

\paragraph{The angels do not know the thoughts of hearts (a.~4).} The
\emph{sed contra} is Jeremias: ``The heart is perverse above all things,
and unsearchable, who can know it? I am the Lord who search the heart, and
prove the reins'' (Jer 17:9--10, Douay); to read the secrets of hearts is
proper to God. A secret thought may be known in two ways. In its effects,
it can be known even by man, and the more subtly as the effect is more
hidden --- by outward act, by change of countenance, by the pulse.
Augustine judges that the demons at times learn men's dispositions, not
only spoken aloud but conceived in thought, when certain signs are
expressed from the soul in the body: \latin{non solum voce prolatas,
verum etiam cogitatione conceptas, cum signa quaedam ex animo exprimuntur
in corpore, tota facilitate perdiscunt} (\emph{De divinatione daemonum}
5.9); Aquinas adds
from the \emph{Retractationes} (II.30, quoted in \emph{De malo} q.~16
a.~8) that Augustine leaves in doubt whether this is by bodily signs
perceptible to them and hidden from us, or by some other spiritual power.
But as they are in the mind, thoughts and the affections of the will God
alone can know: \latin{solus Deus cogitationes cordium et affectiones
voluntatum cognoscere potest}. The reason reaches the heart of the
creature's order: the rational creature is subject to God alone, and He
alone can work within it; ``all that is in the will, and all things that
depend only on the will, are known to God alone,'' and it depends entirely
on the will that a man actually consider what he knows --- ``For what man
knoweth the things of a man, but the spirit of a man that is in him?'' (1
Cor 2:11, Douay). The replies qualify the glory to come. In the
resurrection the body's grossness will no longer hide the soul, and the
brightness of the glorified body will show forth the soul's grace and
glory; Gregory, expounding Job 28:17, says of that society that ``when
the countenance of each one is marked his conscience is penetrated along
with it,'' each distinguishable to another ``as now he cannot be
distinguishable to himself'' (\emph{Moralia} XVIII, \S78, on Job 28:17); yet the
will's closure of its secrets remains even there and even among the
angels, for ``in the resurrection they shall neither marry nor be married,
but shall be as the angels of God in heaven'' (Matt 22:30, Douay) without
therefore reading one another's wills (ad~1). One angel sees another's
intelligible species, but not how far the other actually uses them (ad~2);
and what passes in the sense appetite as it follows bodily impressions the
angels can know, but not as it is moved by the will (ad~3). Damascene's
picture of angelic converse respects the same bound: ``they have no need
of tongue or hearing, but without uttering words they communicate to each
other their own thoughts and counsels'' (\emph{De fide orthodoxa} II.3) ---
they communicate; they do not read.

\paragraph{The angels know the mysteries of grace only by revelation
(a.~5).} The \emph{sed contra} is Dionysius, that Scripture describes
certain heavenly essences ``as questioning Jesus, and learning from Him
the knowledge of His Divine work for us'' (\emph{CH} 7); Parker's text
reads, ``others, as questioning Jesus Himself, as desiring to be
instructed in the science of His Divine work on our behalf, and Jesus
Himself teaching them immediately,'' with Isa 63:1, ``I, that speak
justice'' (cf. Douay). The determination distinguishes the angel's twofold
knowledge. By natural knowledge --- through his essence and through innate
species --- the angels cannot know the mysteries of grace, \latin{mysteria
gratiae}, ``for these mysteries depend upon the pure will of God,'' and
the argument of 1 Cor 2:11 covers God's will as much as man's: ``the
things also that are of God, no man knoweth, but the Spirit of God.''
``There is another knowledge of the angels, which renders them happy''
(\latin{alia Angelorum cognitio, quae eos beatos facit}), the vision of
the Word and of things in the Word; by it they know mysteries of grace,
``but not all mysteries: nor do they all know them equally; but just as
God wills them to learn by revelation,'' the higher beholding the Divine
wisdom more clearly and enlightening the lower. ``Some of these mysteries
they knew from the very beginning of their creation; others they are
taught afterwards, as befits their ministrations.'' So the mystery of the
Incarnation in general was revealed to all from the commencement of their
beatitude, since all their ministry is ordered to it (Heb 1:14), while its
particular conditions even the higher angels learned afterwards (ad~1).
Augustine had already taught the pattern from Eph 3:10, \latin{sic ergo
fuit hoc absconditum a saeculis in Deo, ut tamen innotesceret Principibus
et Potestatibus in coelestibus per Ecclesiam multiformis sapientia Dei} ---
the mystery hidden from ages in God, yet so as to be made known to the
Principalities and Powers in heavenly places through the Church
(\emph{De Genesi ad litteram} V.19.38); the ``great mystery of godliness''
``appeared unto angels'' (1 Tim 3:16, Douay), and the angels desire to
look upon the things now declared in the Gospel (1 Pet 1:12). The blessed
angels behold the Divine wisdom without comprehending it (ad~2); and what
the prophets knew by revelation the angels knew more excellently, the
later prophets more than the earlier (ad~3; Ps 118[119]:100; Aquinas cites
Gregory, \emph{Hom.~xvi in Ezech.}, that ``the knowledge of Divine things
increased as time went on'').

\angeldossier{\ST{I q.~57 aa.~1--5}; parallels \emph{De veritate} q.~8
aa.~11--13; \emph{De malo} q.~16 aa.~7--8; \emph{Summa contra gentiles}
II.99--100.}
 {Common teaching: the angels know material natures (a.~1); they know
 singulars, as their ministry presupposes (a.~2; Heb 1:14; Ps
 90[91]:11); they know the future in its causes but not in itself (a.~3;
 Isa 41:23; Matt 24:36); they do not know the thoughts of hearts (a.~4;
 Jer 17:9--10; 1 Cor 2:11); they know the mysteries of grace only by
 revelation (a.~5).}
 {Ps 90[91]:11 (\emph{sed contra} of a.~2), Heb 1:14, Eccles 5:5, Isa
 41:23 (\emph{sed contra} of a.~3), Matt 24:36, Jer 17:9--10 (\emph{sed
 contra} of a.~4), 1 Cor 2:10--11, Matt 22:30, Isa 63:1, 1 Tim 3:16, Eph
 3:10, 1 Pet 1:12, Amos 3:7, Ps 118[119]:100; Dionysius, \emph{DN} 1
 (a.~1), \emph{CH} 4.2 (a.~1; obj.~3 of a.~5), \emph{CH} 7.3 (\emph{sed
 contra} of a.~5); Augustine, \emph{De divinatione daemonum} 5.9 and
 \emph{Retractationes} II.30 as quoted in \emph{De malo} q.~16 a.~8
 (a.~4), \emph{De Genesi ad litteram} VIII.20.39 (ad~2 of a.~3) and
 V.19.38 (obj.~1 of a.~5); Gregory, \emph{Moralia} XVIII, \S78
 (obj.~1 of a.~4); Damascene, \emph{De fide orthodoxa} II.3.}
 {The unnamed deniers of angelic knowledge of singulars, and Avicenna's
 reduction of that knowledge to universal causes (\emph{De veritate} q.~8
 a.~11 names him), both rejected (a.~2).}

\subsection{The Mode of Angelic Knowledge (\ST{I q.~58})}\label{sec:q58}

Seven articles: whether the angel's intellect is sometimes in potentiality
and sometimes in act; whether he can understand many things at once;
whether his knowledge is discursive; whether he understands by composing
and dividing; whether there can be falsehood in his intellect; and whether
his knowledge is morning and evening knowledge, one or diverse
(\ST{I q.~58 pr.}).

\paragraph{Potentiality and act in the angelic intellect (a.~1).} The
\emph{sed contra} is Augustine: from their creation the angels enjoy the
eternity of the Word in holy and devout contemplation --- in his words,
\latin{ex quo creati sunt, ipsa Verbi aeternitate sancta et pia
contemplatione perfruuntur} (\emph{De Genesi ad litteram} II.8.17). The
intellect is in potentiality in two ways: before it has the habit of
knowledge, and when it has the habit but is not actually considering. In
the first way the angel's intellect is never in potentiality for what his
natural knowledge extends to: as the heavenly bodies have no potentiality
to existence not fully actuated, so ``the heavenly intellects, the angels,
have no intelligible potentiality which is not fully completed by
connatural intelligible species''; yet for things divinely revealed they
may be in potentiality, as even the heavenly bodies are at times in
potentiality to the sun's light. In the second way his intellect can be in
potentiality for what he knows naturally, since he is not always actually
considering all of it; but for the Word and what he beholds in the Word,
never, ``because he is always actually beholding the Word,'' and beatitude
consists in act, not in habit. The angels' desire to look into the
mysteries (1 Pet 1:12) argues no defect, but a delight that excludes
weariness and awaits fresh revelations (ad~2). \emph{Contra Gentiles}
II.97 states the same perfection more absolutely, the separate
substance's intellect being above time: \latin{intellectus substantiae
separatae est semper intelligens actu}.

\paragraph{Many things at once (a.~2).} The \emph{sed contra} is
Augustine, that ``the spiritual faculty of the angelic mind comprehends
most easily at the same time all things that it wills''; his words are,
\latin{secundum potentiam spiritalem mentis angelicae, cuncta quae
voluerit simul notitia facillima comprehendente} (\emph{De Genesi ad
litteram} IV.32). The determination turns on the unity of the object: as
unity of term makes unity of movement, unity of object makes unity of
operation. Distinct things, as distinct, cannot be understood at once; but
in so far as they are comprised under one intelligible species they are
one intelligible object, as the parts of a continuous whole are grasped
together, or subject and predicate in one proposition. Hence by the
knowledge they have of things through the Word, the angels know all things
at once under one species, the Divine essence --- as in heaven, Augustine
says, ``our thoughts will no longer revolve by passing and repassing from
one thing to another, but we shall see all our knowledge at once, and at
one glance'' (\emph{De Trinitate} XV.16.26); by innate species they know
at once all that falls under one species, not what falls under diverse
species.

\paragraph{No discursion (a.~3).} The \emph{sed contra} is Dionysius, that
the angels ``do not acquire Divine knowledge from separate discourses,
nor are they led to something particular from something common''; Parker's
text has them collecting their divine knowledge ``not in portions, or from
portions, or sensible perceptions, or detailed reasonings, or arguing from
something common,'' but contemplating ``intuitively, immaterially and
uniformly'' (\emph{DN} 7.2). The determination perfects the running
analogy of the treatise: as the heavenly bodies have their last perfection
at once from their nature, while earthly bodies reach theirs by change, so
``the lower, namely, the human, intellects obtain their perfection in the
knowledge of truth by a kind of movement and discursive intellectual
operation,'' advancing from one thing known to another; but of the
angels, ``in the truths which they know naturally, they at once behold
all things whatsoever that can be known in them.'' Hence the two names of the two
intellectual natures: the angels ``are called `intellectual beings,''
because what is grasped instantly is said to be understood, and
\latin{intellectus} is the habit of first principles; human souls, which
acquire truth by discursion, are called ``rational,'' \latin{rationales},
``and this comes of the feebleness of their intellectual light.'' The
replies keep the words exact: to know one thing \emph{in} another, as an
object is seen in a mirror, is one motion, not discursion, and so the
angels know creatures in the Word (ad~1); they know syllogisms, and see
effects in causes and causes in effects, without acquiring unknown truth
by syllogizing (ad~2); the ``experience'' Isidore ascribes to demons, as
the objection cites him (\emph{De summo bono} I.10), is said by
similitude, since they know present sensible things ``without any
discursion withal'' (ad~3).

\paragraph{No composing and dividing (a.~4).} The \emph{sed contra} is
Dionysius, that ``the intellectual power of the angel shines forth with
the clear simplicity of divine concepts'' --- Parker: ``they have their
intellectual power and energy re-splendent with the unmixed and undefiled
purity'' (\emph{DN} 7.2). The same weakness of intellectual light that
makes our knowledge discursive makes it compositive: had the intellect, in
apprehending the quiddity of a subject, at once all that can be attributed
to it or removed from it, it would never compose and divide
(\latin{componere et dividere}), but only understand what a thing is. ``In
apprehending a nature, he by one simple perception grasps all that we can
learn by composing and dividing'' (ad~1) --- \latin{intelligendo quod quid
est}. Yet the angel understands our composition and division, as he
apprehends our syllogisms: ``he understands simply, such things as are
composite, things movable immovably, and material things immaterially'';
and when angels speak to men in affirmative and negative propositions,
that shows they know composition and division, not that they know
\emph{by} composing and dividing (ad~3). Dionysius gives the icon of this
intellect when he praises the angel as ``a mirror untarnished----most
transparent----without flaw----pure----without spot'' (\emph{DN} 4.22);
the \emph{Summa} had already spoken the same way of the angelic nature as
``a kind of mirror'' (\latin{quoddam speculum}) representing the Divine
image (q.~56 a.~3).

\paragraph{No falsehood in the angelic intellect (a.~5).} The \emph{sed
contra} joins Aristotle, that ``the intelligence is always true,'' with
Augustine, that nothing can be understood except as it is --- the
\emph{Eighty-Three Questions} read, \latin{non ergo potest quidquam
intellegi nisi ut est} (q.~32). The determination follows from a.~4: the
intellect is always true about what a thing is, as sense about its proper
object; deception creeps in only per accidens, when a quiddity is reached
through composition, as when one definition is taken for another or its
parts do not hang together. Since the angel does not compose and divide,
``no falsehood, error, or deception can exist of itself in the mind of any
angel'' --- \latin{intellectus eius quod quid est semper est verus, nisi
per accidens}. What looks like the contrary among the demons has another
root. ``Owing to their upright will, from their knowing the nature of
every creature, the good angels form no judgments as to the nature of the
qualities therein, save under the Divine ordinance.'' But since ``the
minds of demons are utterly perverted from the Divine wisdom,'' they
judge some things absolutely by their natural conditions; never deceived about natural
properties, they can be misled about supernatural ones --- ``on seeing a
dead man, they may suppose that he will not rise again, or, on beholding
Christ, they may judge Him not to be God.'' Their nescience, which
Dionysius allows the heavenly orders to confess --- ``all the heavenly
Orders are not the same, in all the sacred sciences of the
God-contemplating visions'' (\emph{EH} 6) --- concerns what is knowable
only supernaturally, not their natural knowledge.

\paragraph{Morning and evening knowledge (a.~6).} The \emph{sed contra} is
Augustine, who ``divides the knowledge of the angels into morning and
evening knowledge'' (\emph{De Genesi ad litteram} IV; \emph{De civitate
Dei} XI.7). The expression is Augustine's own device: the six days of
creation are not ordinary days measured by the sun, made only on the
fourth day, but one day --- the day of angelic knowledge directed to the
six classes of things. As morning is the day's beginning and evening its
close, ``their knowledge of the primordial being of things is called
morning knowledge'' (\latin{cognitio matutina}), knowledge of things as
they exist in the Word; ``their knowledge of the very being of the thing
created, as it stands in its own nature, is termed evening knowledge''
(\latin{cognitio vespertina}), because the being of things flows from the
Word as from a primordial principle and is terminated in the being they
have in themselves. Augustine's own sentence fixes the axis: \latin{multum
quippe interest inter cognitionem rei cuiusque in Verbo Dei, et cognitionem
eius in natura eius; ut illud merito ad diem pertineat, hoc ad vesperam} ---
it makes a great difference whether a thing is known in the Word of God
or in its own nature, the former rightly belonging to the day, the latter
to the evening (\emph{De Genesi ad litteram} IV.23.40; the \emph{City of
God} puts it so: ``the knowledge of the creature is, in comparison
of the knowledge of the Creator, but a twilight; and so it dawns and
breaks into morning when the creature is drawn to the praise and love of
the Creator,'' XI.7). The replies answer the anxieties the figure raises.
Evening is not the darkness of error but the end of a day's work; or, as
Augustine allows, one light may be darkness beside a greater --- the
faithful are light beside the wicked (Eph 5:8), yet night beside the life
of glory (2 Pet 1:19) (ad~1). After evening comes no night, because the
good angels, knowing the creature, ``do not adhere to it, for that would
be to turn to darkness and to night; but they refer this back to the
praise of God,'' \latin{post vesperam non ponitur nox, sed mane} --- and
Augustine's own words: \latin{non remanet angelica scientia in eo
quod creatum est, quin hoc continuo referat ad eius laudem atque
caritatem} (\emph{De Genesi ad litteram} IV.24.41); noonday, if one will,
is their knowledge of God Himself (ad~2; \emph{De civitate Dei} XI.29
assigns the angels ``a noonday knowledge'' in God, ``a twilight knowledge''
in themselves). And the third existence of things, in the angelic mind,
counts as evening knowledge, since the angels themselves are creatures
(ad~3).

\paragraph{Whether morning and evening knowledge are one (a.~7).} The
\emph{sed contra} is Augustine: ``There is a vast difference between
knowing anything as it is in the Word of God, and as it is in its own
nature'' (\emph{De Genesi ad litteram} IV). The determination first purges
a misunderstanding: evening knowledge is not drawn \emph{from} the nature
of things --- the angels do not draw knowledge from things at all (q.~55
a.~2) --- but is knowledge of the being things have in their own nature.
And that being the angels know through a twofold medium: by innate
species, and by the forms of things existing in the Word, for in beholding
the Word they see not only the being things have in the Word but the being
things have in themselves, as God by contemplating Himself sees the being
of things. Hence the division: if evening knowledge means knowing the
things' own being \emph{in beholding the Word}, morning and evening
knowledge ``are essentially one and the same, and only differ as to the
things known''; if it means knowing it \emph{through innate ideas}, they
differ --- ``thus Augustine seems to understand it when he assigns one as
inferior to the other.'' The replies dissolve the remaining obstacles:
Gen 1:5, ``there was evening and morning one day,'' takes the day by the
unity of the thing known, which can still be apprehended by diverse ways
of knowing (ad~1); one faculty can have two operations at once when one is
referred to the other, as the will wills end and means together, and the
angels' evening knowledge is referred to their morning knowledge, the
angels always beholding the Father's face (Matt 18:10) (ad~2); and the
perfect's coming does not void this imperfection, for to know a thing in
itself is not opposed to knowing it in its cause, as one conclusion may be
reached by a demonstrative and by a probable medium (ad~3). Augustine had
already refused to put asunder what the angelic mind holds together:
\latin{simul hoc totum possint, simul hoc totum faciant \dots\ simul ergo
habent et diem, et vesperam, et mane} --- they can do all this at once,
and they do; at once they have day, and evening, and morning (\emph{De
Genesi ad litteram} IV.29.46). \emph{De veritate} q.~8 a.~17 adds a
precision the \emph{Summa} leaves implicit: morning and evening, strictly
taken, belong to the day, that is, to the knowledge of grace in the
blessed angels; the demons' natural knowledge, referred to no praise, is
neither.

\angeldossier{\ST{I q.~58 aa.~1--7}; parallels \emph{De veritate} q.~8
aa.~14--17; \emph{Summa contra gentiles} II.97, 101.}
 {Common teaching, after Augustine: the angels' knowledge is morning
 knowledge of things in the Word and evening knowledge of things in their
 own nature (aa.~6--7; \emph{De Genesi ad litteram} IV; \emph{De civitate
 Dei} XI.7, 29). Thomist position: the angelic intellect is in
 potentiality only as not always considering what it knows (a.~1); it
 understands many things at once under one species (a.~2); it knows
 without discursion (a.~3), without composing and dividing (a.~4), and
 without the possibility of falsehood of itself (a.~5).}
 {Augustine, \emph{De Genesi ad litteram} II.8.17 (\emph{sed contra} of
 a.~1), IV.32 (a.~2), IV.23--24 (aa.~6--7), IV.29.46 (a.~7), \emph{De
 Trinitate} XV.16.26 (a.~2), \emph{De civitate Dei} XI.7 and XI.29 (a.~6),
 \emph{De diversis quaestionibus octoginta tribus} q.~32 (\emph{sed
 contra} of a.~5); Dionysius, \emph{DN} 7.2 (\emph{sed contra} of aa.~3
 and~4), \emph{DN} 4.8 (obj.~1 of a.~1) and 4.22 (a.~4), \emph{DN} 4
 (objs.~1 and~3 of a.~5), \emph{EH} 6 (obj.~2 of a.~5); 1 Pet 1:12, Gen
 1:5, Matt 18:10, Eph 5:8, 2 Pet 1:19, 1 Cor 13:10; Aristotle and Isidore
 as cited by Aquinas.}
 {Isidore's ``experience'' in demons, received only as a similitude
 (a.~3); the demons' accidental error about supernatural matters (a.~5);
 \emph{De veritate} q.~8 a.~17 restricts morning and evening knowledge,
 strictly taken, to the blessed (a.~7).}
